\cotefile{148}
\scanpage{11}
\mspage{4}

Le groupe $\mathbb Z/6\mathbb Z$ opère transitivement sur $({}_2X_S)\setminus\{0\}$ car $\pm1$ opère trivialt - donc ${}_2X_S\setminus\{0\}$ est un torseur sous
\[
(\mathbb Z/6\mathbb Z)/\pm1\simeq\mathbb Z/3\mathbb Z
\]
\deleted{\unsure{(conj.|conj.}} (générateur correspondant $\unsure{\simeq|=}\,-\bar\zeta\,\unsure{j|\zeta}=\exp\frac{2\pi}{6}$),
\deleted{\unsure{en ce cas} torseur} \inserted{le torseur} s'identifie : $S/\pm1$, en
% Source-state: two short crossed strokes beneath a raised replacement a;
% the old strokes are reproduced without assigning them a letter value.
\begin{tikzpicture}[baseline={(0,0)},x=1em,y=1ex,line width=.35pt]
  \draw (.08,2.25) .. controls (.16,2.33) and (.25,2.32) .. (.34,2.28);
  \draw (.09,1.78) -- (.34,1.80);
  \draw (-.02,1.88) -- (1.34,2.31);
  \node[anchor=base,inner sep=0pt] at (.99,2.08) {\inserted{a}};
\end{tikzpicture}ssociant \unsure{:|à} chaque $s\in S$ le milieu du rayon $[0,s]$, ou plutôt son image dans\inserted{)} $X_S$.

L'ens des pts fixes de $\mathbb Z/3\mathbb Z\inserted{\subset\mathbb Z/6\mathbb Z}$ sur $X_S$ est de card 3, il s'obtient en prenant les \unsure{barycentres} des six triangles de la subdivision centrale (bibarycentrique) de l'hexagone (pts cerclés) (tels que $(1-\bar\zeta)/3$ etc dans le cas épinglé), ils se groupent en deux paquets de 3 \unsure{mutuellement|mutuellem\textsuperscript{t}} congrus sous $\Gamma$, \unsure{donnant|donnent} deux éléments (symétriques) de $X_S$, d'ordre 3,

\margin{10}
